% Darstellungen des begrenzten Cantor--Bernstein-Piloten.
% Die regulären Fachbände enthalten Aussagen und Verweise auf die Begleitbeweise.
\newif\ifCSBReferenz
\newif\ifCSBLesetext
\newif\ifCSBTabellen
\CSBReferenztrue
\CSBLesetextfalse
\CSBTabellenfalse
\NewDocumentCommand{\CSBProof}{+m}{\ifCSBTabellen#1\fi}
% Textüberschriften ohne neue Hyperref-Zähler: bestehende Satzziele im
% integrierten Band bleiben auch nach Einfügen des Lesetexts unverändert.
\newcommand{\CSBHeading}[1]{\par\Needspace{5\baselineskip}\medskip\noindent\textbf{#1}\par\smallskip}

% Pfade gelten zunächst für die Build-Produkte. publish-pdfs.py ersetzt sie
% bei der Veröffentlichung durch die lesbaren benachbarten PDF-Dateinamen.
\begin{luacode*}
  local job = tex.jobname or ""
  local directory = "registry/csb/"
  if job:match("^_B[0-9][0-9]$") then directory = "csb/" end
  if job:match("^_B08%-csb%-") then directory = "" end
  tex.sprint("\\gdef\\CSBEditionDirectory{" .. directory .. "}")
\end{luacode*}
\newcommand{\CSBReadingLink}{%
  \href{\CSBEditionDirectory_B08-csb-reading.pdf\#csb.reading}{Lesefassung}%
}
\newcommand{\CSBProofsLink}{%
  \href{\CSBEditionDirectory_B08-csb-proofs.pdf\#csb.proofs}{Beweistabellen}%
}

% Die drei vorhandenen H-Identitäten bleiben im Fachband als Beweisverweise
% erreichbar. Die eigentlichen Hilfsbeweise stehen nur in der Tabellenfassung.
\newcommand{\CSBAuxReference}[1]{%
  \refstepcounter{proofpartnr}%
  \edef\mxPPLbl{thm:pp:\mxProofPartParent:\arabic{proofpartnr}}%
  \label{\mxPPLbl}%
  \ThmLookupRegisterId{#1}{theorem}{\mxPPLbl}%
  \noindent Hilfsresultat \ThmLookupEmitRef{theorem}{\mxPPLbl}: %
  \href{\CSBEditionDirectory_B08-csb-proofs.pdf\#csb.helper.#1}{%
    Aussage und Beweis in den Beweistabellen}.\par
}
